% Part: history
% Chapter: biographies
% Section: emmy-noether

\documentclass[../../../include/open-logic-section]{subfiles}

\begin{document}

\olfileid{his}{bio}{noe}

\olsection{إيمي نويثر}

\olphoto{noether-emmy}{Emmy Noether}

وُلدت إيمي نويثر (\textsc{ner}-ter) بإرلنغن من بلاد ألمانيا في 23 مارس
1882، ونشأت في بيت من بيوت العلم ميسور الحال. وقد استحقت أن تُلقَّب «أم
الجبر الحديث»، لما أحدثته من آثار باقية في الرياضيات والفيزياء جميعاً،
على كثرة ما كان يعترض النساء يومئذ من عوائق التعلّم. إذ جرت عادة الألمان
حينئذ أن يُقصَر تعليم الفتيات على الفنون، وأن يُحلن دون المدارس المهيئة
للجامعة. غير أن نويثر كانت تحضر الدروس مستمعة في جامعتي غوتنغن وإرلنغن،
وكان أبوها أستاذاً للرياضيات في الثانية؛ ثم لما غيّرت الجامعة سياستها
وأذنت بقبول النساء، قُيّدت فيها طالبة سنة 1904، ونالت الدكتوراه في
الرياضيات سنة 1907.

ولم تشفع لها أهليتها عند من قاوموها في أثناء عملها؛ فكانت تدرّس بإرلنغن
من سنة 1908 إلى سنة 1915 ولا رزق لها على ذلك. وفي هذه المدة عرف فضلها
ديفيد هيلبرت، وكان من أعلام الرياضيين في عصره، فاستدعاها إلى غوتنغن. لكن
منع النساء من رتبة الأستاذية ألجأها إلى أن تلقي الدروس باسم هيلبرت، من غير
أجر أيضاً. وفي تلك الأيام أثبتت القضية التي تُعرف الآن بمبرهنة نويثر، وما
برحت أصلاً معمولاً به في الفيزياء النظرية. ثم أُذن لها بالتدريس سنة 1919.
ويُحكى أن هيلبرت، لما رأى زملاءه لا يزالون على اعتراضهم، قال: «أيها
السادة، إن مجلس الكلية ليس حمّاماً عاماً».

ثم صرفت عنايتها، في أواخر عشرينيات القرن العشرين، إلى الجبر المجرد، فقلبت
آثارها وجه هذا العلم. وكان من الأدوات التي يكثر دورانها في براهينها ما
يسمى شرط السلسلة الصاعدة: ألا توجد سلسلة لا نهاية لها من بعض المجموعات،
كل واحد منها داخل فيما بعده دخولا صارماً. ومن ذلك أن البنى التي تسمى الآن
حلقات نويثرية لا تقبل متتاليات لا نهائية من المثاليات
$I_1 \subsetneq I_2 \subsetneq \dots$. وليس هذا الشرط مقصوراً على ذلك،
بل يعم كل ترتيب جزئي؛ وإنما صورته الجبرية الخاصة ترتيب المثاليات بعلاقة
الاحتواء. ويقابله شرط السلسلة النازلة، وهو أن تبلغ كل متتالية
\emph{متناقصة} بصرامة في ترتيب جزئي منتهاها بعد عدد من الخطوات. فإذا وفّى
ترتيب جزئي بهذا الشرط، جاز إجراء الاستقراء فيه على نحو يشبه إجراءنا إياه
في الترتيب $<$ على~$\Nat$. وتسمى أمثال هذه الترتيبات \emph{حسنة
التأسيس} أو \emph{نويثرية}، ويسمى أصل البرهان المقابل \emph{الاستقراء
النويثري}.

وكانت نويثر يهودية؛ فلما استولى النازيون على الحكم سنة 1933 عُزلت من
منصبها. ثم تيسّر لها أن تهاجر إلى الولايات المتحدة، فتولت عملاً مؤقتاً في
كلية برين ماور بولاية بنسلفانيا. وكانت، وهي مقيمة هناك، تحاضر في برينستون
أيضاً، وإن وجدت جامعتها نافرة من النساء \citep[81]{Dick1981}. وفي سنة
1935 أُجريت لها جراحة لاستئصال ورم في الرحم، فأصابتها عدوى من أثرها،
وماتت، ودُفنت في برين ماور.

\begin{reading}
ومن أراد سيرة نويثر فلينظر \citet{Dick1981}. وقد أتاح معهد بيريميتر
للفيزياء النظرية على الشبكة محاضراته في حياتها وأثرها
\citep{Perimeter2015}. ومن آثر السماع على القراءة وجد في \emph{أشياء
فاتتك في حصة التاريخ} تدوينة صوتية تتناول حياتها وأثرها
\citep{Frey2015}. ومجموع آثار نويثر منشور بلغته الألمانية الأصلية
\citep{Noether1983}.
\end{reading}

\end{document}
